\chapter{Valeurs manquantes}
\label{ch:missing-data}

\begin{quote}
\emph{<< Les donn\'ees que vous n'avez pas sont souvent plus importantes que celles que vous avez. >>}
--- Statisticien anonyme
\end{quote}

% --- hook ---
En 1976, Donald Rubin publie un papier qui devrait être lu par toute personne qui touche à un jeu de données. L'idée est simple : une valeur manquante n'est pas du bruit, c'est de l'information --- et l'information la plus importante n'est pas la valeur qu'elle aurait dû prendre, c'est la \emph{raison} pour laquelle elle manque. Rubin formalise trois mécanismes --- MCAR, MAR, MNAR --- et démontre que l'imputation la plus naturelle (remplacer par la moyenne) est valide pour le premier, biaisée pour le deuxième, et catastrophique pour le troisième. Cinquante ans plus tard, la moitié des modèles de machine learning ignorent encore cette distinction. Ce chapitre la traite avec le sérieux qu'elle mérite.
\section{Pourquoi des donn\'ees manquent}

Une valeur manquante n'est pas une nuisance --- c'est une information. La \emph{raison} pour laquelle une donn\'ee manque d\'etermine comment vous devriez la traiter. Trois m\'ecanismes ont \'et\'e formalis\'es par Rubin (1976) :

\begin{definition}[M\'ecanismes de donn\'ees manquantes]
\begin{itemize}
  \item \textbf{MCAR} (Missing Completely At Random) : la manquance est ind\'ependante de toute variable, observ\'ee ou non. Exemple : un \'echantillon de laboratoire est accidentellement tomb\'e.
  \item \textbf{MAR} (Missing At Random) : la manquance d\'epend de variables observ\'ees mais pas de la valeur manquante elle-m\^eme. Exemple : les patients plus jeunes sautent un questionnaire de d\'epression.
  \item \textbf{MNAR} (Missing Not At Random) : la manquance d\'epend de la valeur non observ\'ee. Exemple : les hauts revenus refusent de d\'eclarer leur revenu.
\end{itemize}
\end{definition}

\begin{warningbox}
On ne peut pas prouver statistiquement le MNAR --- cela requiert de la connaissance m\'etier. Mais la distinction compte \'enorm\'ement : l'imputation simple marche pour MCAR, est raisonnable pour MAR, et peut \^etre fortement biais\'ee pour MNAR.
\end{warningbox}

% ============================================================
\section{D\'etecter les valeurs manquantes}

\begin{minted}{python}
import pandas as pd
import numpy as np

# Charger les donn\'ees Titanic
url = ("https://raw.githubusercontent.com/datasciencedojo/"
       "datasets/master/titanic.csv")
df = pd.read_csv(url)

# Compter les valeurs manquantes
print(df.isnull().sum())
print()
print(f"Total missing cells: {df.isnull().sum().sum()}")
print(f"Percentage missing: {df.isnull().mean().mean() * 100:.1f}%")
\end{minted}

\begin{minted}{python}
# Rapport d\'etaill\'e des manquantes
def missing_report(df):
    """Return a DataFrame summarizing missing data."""
    missing = df.isnull().sum()
    pct = (missing / len(df)) * 100
    report = pd.DataFrame({"missing": missing, "pct": pct.round(1)})
    return report[report["missing"] > 0].sort_values("pct", ascending=False)

print(missing_report(df))
\end{minted}

\subsection{Manquantes cach\'ees}

\begin{minted}{python}
# Les valeurs manquantes ne sont pas toujours NaN
# D\'eguisements courants : "?", "N/A", "", -1, 999, -999, "unknown"
df_adult = pd.read_csv(
    "https://archive.ics.uci.edu/ml/machine-learning-databases/adult/adult.data",
    header=None,
    na_values=["?", " ?"],  # Adult utilise " ?" pour les manquantes
    names=["age", "workclass", "fnlwgt", "education", "education_num",
           "marital_status", "occupation", "relationship", "race",
           "sex", "capital_gain", "capital_loss", "hours_per_week",
           "native_country", "income"]
)
print(missing_report(df_adult))
\end{minted}

\begin{datatip}
Scrutez toujours les valeurs sentinelles : \texttt{-1}, \texttt{999}, \texttt{0} dans les colonnes num\'eriques ; cha\^ines vides, << unknown >>, << N/A >> dans les colonnes texte. Remplacez-les par \texttt{np.nan} avant toute analyse.
\end{datatip}

% ============================================================
\section{Visualiser les manquantes}

\begin{minted}{python}
import missingno as msno
import matplotlib.pyplot as plt

# Matrix plot : lignes blanches = manquantes
msno.matrix(df, figsize=(10, 5), sparkline=False)
plt.title("Missing data pattern --- Titanic")
plt.tight_layout()
plt.savefig("missing_matrix.png", dpi=150)
plt.show()
\end{minted}

\begin{minted}{python}
# Bar chart : compte de non-nuls par colonne
msno.bar(df, figsize=(10, 4))
plt.tight_layout()
plt.show()

# Heatmap : corr\'elation de manquance entre colonnes
# Valeurs proches de +1 = colonnes manquantes ensemble
msno.heatmap(df, figsize=(8, 6))
plt.tight_layout()
plt.show()
\end{minted}

\begin{minted}{python}
# Dendrogramme : clustering hi\'erarchique des patterns de manquance
msno.dendrogram(df, figsize=(10, 5))
plt.tight_layout()
plt.show()
\end{minted}

\begin{preprocesstip}
La heatmap de corr\'elation de manquance est extr\^emement utile : si deux colonnes ont des manquances corr\'el\'ees, elles partagent probablement une cause commune (par ex. m\^eme section de questionnaire que certains r\'epondants ont saut\'ee).
\end{preprocesstip}

% ============================================================
\section{Strat\'egies de suppression}

\subsection{Suppression listwise (lignes)}

\begin{minted}{python}
# Supprimer les lignes avec UNE valeur manquante
df_complete = df.dropna()
print(f"Before: {len(df)}, After: {len(df_complete)}")
print(f"Lost: {len(df) - len(df_complete)} rows "
      f"({(len(df) - len(df_complete)) / len(df) * 100:.1f}%)")
\end{minted}

\subsection{Suppression de colonnes}

\begin{minted}{python}
# Supprimer les colonnes \`a plus de 50% manquantes
threshold = 0.5
cols_to_drop = df.columns[df.isnull().mean() > threshold].tolist()
print(f"Dropping columns: {cols_to_drop}")
df_reduced = df.drop(columns=cols_to_drop)
\end{minted}

\begin{warningbox}
La suppression listwise n'est s\^ure que sous MCAR. Sous MAR ou MNAR, elle introduit un biais. Pour Titanic, les passagers plus \^ag\'es en classes sup\'erieures ont plus souvent un \^age renseign\'e --- supprimer les \^ages manquants biaise le dataset vers les classes inf\'erieures.
\end{warningbox}

% ============================================================
\section{Imputation simple}

\subsection{Imputation par moyenne, m\'ediane, mode}

\begin{minted}{python}
from sklearn.impute import SimpleImputer

# Num\'erique : imputation m\'ediane (robuste aux outliers)
num_imputer = SimpleImputer(strategy="median")
df["Age"] = num_imputer.fit_transform(df[["Age"]])

# Cat\'egoriel : imputation par mode (le plus fr\'equent)
cat_imputer = SimpleImputer(strategy="most_frequent")
df["Embarked"] = cat_imputer.fit_transform(df[["Embarked"]]).ravel()

print(df.isnull().sum())
\end{minted}

\subsection{Imputation constante}

\begin{minted}{python}
# Remplir avec une valeur sp\'ecifique
df["Cabin"] = df["Cabin"].fillna("Unknown")

# Remplir num\'erique avec 0 (\`a utiliser avec prudence)
df["some_col"] = df["some_col"].fillna(0)
\end{minted}

% ============================================================
\section{Imputation avanc\'ee}

\subsection{Imputation KNN}

\begin{minted}{python}
from sklearn.impute import KNNImputer

# KNN utilise des lignes similaires pour estimer les manquantes
# Ne fonctionne que sur des donn\'ees num\'eriques
num_cols = ["Age", "Fare", "SibSp", "Parch"]
knn_imputer = KNNImputer(n_neighbors=5, weights="distance")
df[num_cols] = knn_imputer.fit_transform(df[num_cols])

print(f"Missing after KNN: {df[num_cols].isnull().sum().sum()}")
\end{minted}

\subsection{Imputation it\'erative (MICE)}

\begin{minted}{python}
from sklearn.experimental import enable_iterative_imputer
from sklearn.impute import IterativeImputer

# L'imputeur it\'eratif mod\'elise chaque variable en fonction des autres
iter_imputer = IterativeImputer(max_iter=10, random_state=42)
df[num_cols] = iter_imputer.fit_transform(df[num_cols])

print("Iterative imputation complete.")
print(df[num_cols].describe())
\end{minted}

\begin{preprocesstip}
L'imputation KNN respecte la structure locale (des passagers similaires re\c{c}oivent des \^ages similaires). L'imputation it\'erative capture les relations lin\'eaires entre variables. Les deux surpassent l'imputation moyenne/m\'ediane quand les donn\'ees sont MAR.
\end{preprocesstip}

% ============================================================
\section{Imputation avec indicateur}

\begin{minted}{python}
# Cr\'eer un indicateur binaire AVANT d'imputer
df["Age_was_missing"] = df["Age"].isnull().astype(int)

# Puis imputer
df["Age"] = df["Age"].fillna(df["Age"].median())

# L'indicateur pr\'eserve l'information que la valeur \'etait manquante
print(df[["Age", "Age_was_missing"]].head(10))
\end{minted}

\begin{datatip}
Ajouter une colonne indicatrice est l'une des techniques les plus sous-utilis\'ees. Elle laisse les mod\`eles en aval apprendre si la manquance elle-m\^eme est pr\'edictive. Pour Titanic, avoir un num\'ero de cabine manquant est corr\'el\'e \`a une moindre survie.
\end{datatip}

% ============================================================
\section{Choisir une strat\'egie d'imputation}

\begin{longtable}{p{3.5cm}p{3cm}p{3cm}p{3.5cm}}
\toprule
\textbf{M\'ethode} & \textbf{Bon pour} & \textbf{Faiblesse} & \textbf{M\'ecanisme} \\
\midrule
Supprimer lignes & $<$5\% manquantes & Perte de donn\'ees & MCAR seul \\
Moyenne/m\'ediane & Baseline rapide & R\'eduit la variance & MCAR \\
Mode & Cat\'egoriel & Ignore la structure & MCAR \\
KNN & Donn\'ees clusteris\'ees & Lent sur grand volume & MAR \\
It\'eratif (MICE) & Relations lin\'eaires & Suppose lin\'earit\'e & MAR \\
Constante m\'etier & Sens connu & Peut biaiser & Tout \\
\bottomrule
\end{longtable}

% ============================================================
\section{Exercices}

\begin{exercise}
\begin{enumerate}
  \item Chargez le dataset Melbourne Housing. Cr\'eez un rapport de manquantes et un missingno matrix plot. Quelles colonnes ont des manquances corr\'el\'ees ?
  \item Pour Titanic, comparez la distribution de \texttt{Age} avant et apr\`es imputation moyenne vs imputation KNN. Tracez les deux distributions sur le m\^eme histogramme.
  \item Chargez Adult Census (avec \texttt{na\_values=[" ?"]}) et imputez les trois colonnes cat\'egorielles \`a valeurs manquantes par mode. V\'erifiez qu'il ne reste aucune manquante.
  \item Impl\'ementez une fonction \texttt{smart\_impute(df)} qui : (a) supprime les colonnes \`a $>$70\% manquantes, (b) impute les num\'eriques par m\'ediane, (c) impute les cat\'egorielles par mode, (d) ajoute des indicateurs pour toutes les colonnes initialement manquantes.
  \item Investiguez si \texttt{Age} dans Titanic est MCAR en comparant le taux de survie des passagers avec et sans \^age renseign\'e. Que concluez-vous ?
\end{enumerate}
\end{exercise}

% ============================================================
\section{R\'esum\'e du chapitre}

\begin{itemize}
  \item Les donn\'ees manquantes ont trois m\'ecanismes : MCAR, MAR et MNAR. Le m\'ecanisme d\'etermine quelle imputation est valide.
  \item Scrutez toujours les valeurs sentinelles cach\'ees (\texttt{?}, \texttt{-1}, \texttt{999}).
  \item La biblioth\`eque \texttt{missingno} fournit visualisations matrix, bar, heatmap et dendrogramme.
  \item L'imputation simple (moyenne, m\'ediane, mode) est rapide mais r\'eduit la variance et ne marche bien que sous MCAR.
  \item KNN et MICE capturent les relations entre variables et fonctionnent sous MAR.
  \item Ajouter une colonne indicatrice de manquance pr\'eserve l'information contenue dans le fait qu'une valeur \'etait manquante.
\end{itemize}
